\begin{ouverture}[courbe]{1}{Nombres rationnels et fractions}{Partie 1 · Nombres et calculs}
{Signes · Priorités · Fractions · Calcul exact}
\objectif{Calculer en respectant les priorités.}
\objectif{Choisir entre valeur exacte et valeur approchée.}
\objectif{Résoudre un problème avec des fractions.}
\prerequis{Tables de multiplication et nombres relatifs.}
\end{ouverture}
\pagestyle{otMain}\essentiel\fontsize{11}{15}\selectfont
\essSection{1}{Organiser un calcul}
\begin{definition}Un nombre rationnel peut s’écrire $a/b$, où $a$ et $b$ sont des entiers et $b\ne0$. Deux fractions représentent le même nombre si on multiplie ou divise leurs deux termes par un même nombre non nul.\end{definition}
\begin{propriete}Les parenthèses sont prioritaires, puis les puissances, les produits et quotients, enfin les sommes et différences. À priorité égale, on calcule de gauche à droite.\end{propriete}
\begin{exemple}$-4+3\times(5-7)=-4+3\times(-2)=-10$. Pour $(-3)^2$, le carré porte sur le nombre $-3$ ; dans $-3^2$, le signe moins est extérieur au carré.\end{exemple}
\begin{savoirfaire}[Gérer les signes]Un produit ou un quotient de deux nombres de même signe est positif ; de signes contraires, il est négatif. Pour une somme, ces règles ne s’appliquent pas : comparer les valeurs et les signes.\end{savoirfaire}
\begin{attention}On ne divise jamais par zéro. Une égalité doit relier des nombres égaux : ne pas enchaîner un résultat partiel avec tout le calcul initial.\end{attention}
\clearpage
\essSection{2}{Calculer avec des fractions}
\begin{propriete}Pour additionner, utiliser un dénominateur commun. Pour multiplier, multiplier les numérateurs et les dénominateurs. Pour diviser par une fraction non nulle, multiplier par son inverse.\end{propriete}
\begin{exemple}
\[\frac34-\frac56=\frac9{12}-\frac{10}{12}=-\frac1{12},
\qquad \frac23\div\frac57=\frac23\times\frac75=\frac{14}{15}.\]
\end{exemple}
\begin{savoirfaire}[Une fraction d’une quantité]Prendre $3/5$ de 40 revient à calculer $40\times3/5=24$. « Les $2/3$ du reste » impose de calculer d’abord le reste ; le tout de référence a changé.\end{savoirfaire}
\begin{propriete}Comparer des fractions de dénominateurs positifs en les réduisant au même dénominateur. Une valeur exacte comme $7/3$ se distingue de son arrondi $2{,}33$.\end{propriete}
\begin{attention}On simplifie des facteurs dans un produit, jamais des termes isolés d’une somme. Ainsi $(6+3)/3=3$, et non $6+1$.\end{attention}
\otIntertitre{Contrôler sa réponse}
Estimer le signe et l’ordre de grandeur avant le calcul. Une fraction d’une quantité positive, comprise entre 0 et 1, donne une quantité plus petite que le tout.
\clearpage\rubrique{Découverte}
\begin{blocLivre}{ex-1-1}\begin{exercice}[1]{Une température}
Le matin, il fait $-6\,^\circ$C. La température augmente de 11 degrés puis baisse de 4 degrés. Quelle est la température finale ?
\end{exercice}\end{blocLivre}
\begin{blocLivre}{ex-1-2}\begin{exercice}[1]{Respecter les priorités}
Calculer $A=18-3(2+5)$ et $B=(18-3)(2+5)$.
\end{exercice}\end{blocLivre}
\begin{blocLivre}{ex-1-3}\begin{exercice}[1]{Deux écritures, un nombre}
Compléter $\dfrac5{8}=\dfrac{\,?\,}{40}$ et $\dfrac{42}{63}=\dfrac2{\,?\,}$.
\end{exercice}\end{blocLivre}
\begin{blocLivre}{ex-1-4}\begin{exercice}[1]{Une somme exacte}
Calculer $\dfrac7{12}+\dfrac5{18}$ et donner une fraction simplifiée.
\end{exercice}\end{blocLivre}
\begin{blocLivre}{ex-1-5}\begin{exercice}[1]{Produit et quotient}
Calculer $\dfrac{-9}{14}\times\dfrac7{15}$ puis $\dfrac35\div\dfrac{-9}{10}$.
\end{exercice}\end{blocLivre}

\clearpage\rubrique{Entraînement}
\begin{blocLivre}{ex-1-6}\begin{exercice}[2]{Comparer sans arrondir}
Ranger $-\dfrac34$, $-\dfrac56$ et $-\dfrac7{12}$ dans l’ordre croissant.
\end{exercice}\end{blocLivre}
\begin{blocLivre}{ex-1-7}\begin{exercice}[2]{Une expression à étages}
Calculer exactement $C=\dfrac{\frac23+\frac16}{\frac54-\frac12}$.
\end{exercice}\end{blocLivre}
\begin{blocLivre}{ex-1-8}\begin{exercice}[2]{Une réduction trompeuse}
Un élève écrit $(8+4)/4=8+1$. Expliquer son erreur et calculer correctement.
\end{exercice}\end{blocLivre}
\begin{blocLivre}{ex-1-9}\begin{exercice}[2]{Une recette}
Une recette pour 6 personnes utilise $3/4$ L de lait. Quel volume faut-il pour 10 personnes ?
\end{exercice}\end{blocLivre}
\begin{blocLivre}{ex-1-10}\begin{exercice}[2]{Exact ou approché ?}
Donner la valeur exacte de $2-\dfrac57$, puis son arrondi au centième.
\end{exercice}\end{blocLivre}

\clearpage\rubrique{Problèmes}
\begin{blocLivre}{ex-1-11}\begin{exercice}[3]{Deux étapes de lecture}
Un lecteur lit $2/5$ d’un livre le lundi, puis $1/3$ de ce qui reste le mardi. Quelle fraction du livre reste-t-il ? Combien de pages pour un livre de 270 pages ?
\end{exercice}\end{blocLivre}
\begin{blocLivre}{ex-1-12}\begin{exercice}[3]{Une réserve d’eau}
Une cuve est remplie aux $7/8$. On consomme $2/7$ de l’eau présente, puis on ajoute 30 L. La cuve est alors remplie aux $3/4$. Quelle est sa capacité ?
\end{exercice}\end{blocLivre}
\begin{blocLivre}{ex-1-13}\begin{exercice}[3]{Partager un trajet}
Un cycliste parcourt $3/8$ du trajet le matin et $5/12$ l’après-midi. Les 15 km restants sont faits le soir. Calculer la longueur totale.
\end{exercice}\end{blocLivre}
\begin{blocLivre}{ex-1-14}\begin{exercice}[3]{Des votes et des absents}
Dans une classe, $5/6$ des élèves sont présents. Parmi eux, $3/5$ choisissent le projet A. Le projet A reçoit 15 voix. Combien y a-t-il d’élèves et d’absents ?
\end{exercice}\end{blocLivre}
\begin{blocLivre}{ex-1-15}\begin{exercice}[3]{Trouver toutes les possibilités}
On cherche les entiers naturels $n$ tels que $n/12$ soit strictement supérieur à $2/3$ et strictement inférieur à $5/4$. Les donner tous.
\end{exercice}\end{blocLivre}

\clearpage
